\documentclass[11pt]{amsart}

\usepackage[margin=1.15in]{geometry}
\usepackage{amsmath,amssymb,amsthm}
\usepackage{mathtools}
\usepackage{booktabs}
\usepackage[hidelinks]{hyperref}
\usepackage{xcolor}
\emergencystretch=2em

\newtheorem{theorem}{Theorem}[section]
\newtheorem{lemma}[theorem]{Lemma}
\newtheorem{proposition}[theorem]{Proposition}
\newtheorem{corollary}[theorem]{Corollary}
\theoremstyle{definition}
\newtheorem{definition}[theorem]{Definition}
\theoremstyle{remark}
\newtheorem{remark}[theorem]{Remark}

\newcommand{\lean}[1]{\texttt{#1}}
\DeclareMathOperator{\Good}{Good}
\DeclareMathOperator{\Bad}{Bad}

\title[A discrete iteration related to Pierce expansions]{The iteration $x\mapsto a \bmod x$ has length $O_\varepsilon(a^{3/10+\varepsilon})$}

\author{Jason Davies}
\email{jason@jasondavies.com}

\date{October 2026}

\begin{document}

\begin{abstract}
For integers $a>b>0$ let $b_1=b$, $b_{i+1}=a\bmod b_i$, and let $P(a,b)$ be the number of
positive terms of this sequence. Shallit asked whether $P(a,b)=O((\log a)^2)$. The bound
$P(a,b)\ll_\varepsilon a^{1/3+\varepsilon}$ of Erd\H{o}s and Shallit was improved by Chase and
Pandey to $a^{1/3-2/177+\varepsilon}$. We prove $P(a,b)\ll_\varepsilon a^{3/10+\varepsilon}$
by an elementary argument. The new ingredient is an exact identity for three consecutive steps
of the iteration. It shows that a certain integer combination $D=dv-eu$ of the drops and
quotient increments is always positive. It is also small whenever the three drops are small,
and $D$ together with the two increments confines the block to a curve with few integer points.
The whole proof is formalised in Lean~4 with Mathlib, without \lean{sorry}.
\end{abstract}

\maketitle

\section{Introduction}

Let $a>b>0$ be integers. Put $b_1=b$ and $b_{i+1}=a \bmod b_i$ while $b_i>0$. The sequence is
strictly decreasing until it reaches $0$; if $b_{n+1}=0$ we write $P(a,b)=n$. For example
$P(35,22)=7$, from $22,13,9,8,3,2,1,0$. The number $P(a,b)$ is the length of the Pierce
expansion of $b/a$~\cite{Pierce,Shallit86}, since
\[
\frac ba=\frac1{q_1}-\frac1{q_1q_2}+\frac1{q_1q_2q_3}-\cdots,\qquad q_i=\lfloor a/b_i\rfloor .
\]
Shallit~\cite{OPG} asked whether
\begin{equation}\label{eq:conj}
P(a,b)=O\bigl((\log a)^2\bigr).
\end{equation}

Shallit~\cite{Shallit86} proved $P(a,b)\ll a^{1/2}$ by ``Archimedean'' arguments.
Erd\H{o}s and Shallit~\cite{ES} combined these with a divisor argument to get
$P(a,b)\ll_\varepsilon a^{1/3+\varepsilon}$, and showed that $\max_b P(a,b)/\log a$ does not
tend to $0$. Chase and Pandey~\cite{CP} went past the exponent $1/3$, proving
$P(a,b)\ll_\varepsilon a^{1/3-2/177+\varepsilon}$ with exponential sum estimates
(Bourgain's exponent pair), and showed that $\max_b P(a,b)\gg \log a/\log\log a$ for all large
$a$. Shallit's talk pages~\cite{ShallitTalks} report a bound $n^{3/10}$, described as
``LLM, unpublished so far''. We have not seen that argument. The result below has the same
exponent up to $\varepsilon$, and we make no claim about how the two arguments are related.

\begin{theorem}\label{thm:main}
For every $\varepsilon>0$ there is a constant $C_\varepsilon$ such that
\[
P(a,b)\le C_\varepsilon\, a^{3/10+\varepsilon}\qquad\text{for all integers }0<b<a .
\]
\end{theorem}

The proof is elementary. It uses only the division equations
$a=q_ib_i+b_{i+1}$ and the divisor bound $\tau(n)\ll_\delta n^\delta$. As usual, the remainders
are split into dyadic windows $[A,2A)$, and we bound the number $T$ of remainders in one
window. There are two bounds:
\begin{align}
T(T-1)&\le 4A\,\Delta, \label{eq:intro1}\\
T^6A^2&\le 995328\, a^3\Delta \qquad(T\ge 12), \label{eq:intro2}
\end{align}
where $\Delta$ is the largest value of the divisor function on $[1,32a^7]$.
Bound~\eqref{eq:intro1} is the Erd\H{o}s--Shallit argument. Bound~\eqref{eq:intro2} is new.
Multiplying~\eqref{eq:intro2} by the square of~\eqref{eq:intro1} eliminates $A$ and gives
$T^{10}\ll a^3\Delta^3$, which is the exponent $3/10$.

The new bound comes from three consecutive steps $x=b_i$, $y=b_{i+1}$, $z=b_{i+2}$,
$w=b_{i+3}$. Let $d,e,f$ be the drops, let $u,v$ be the increments of the quotients, and put
$D=dv-eu$. Then
\begin{equation}\label{eq:intro-curv}
D\,z=df+ude+(u-1)e^2 .
\end{equation}
So $D$ is a \emph{positive} integer, even though $dv$ and $eu$ largely cancel, and it is small
when the drops are small. Also $(q,x)=(\lfloor a/b_i\rfloor,b_i)$ lies on the curve
$xM(q)=aN(q)+D$, where $M$ and $N$ are explicit polynomials depending on $(u,v)$
(Lemma~\ref{lem:fiber-eq}). A resultant-type identity shows that each such curve has
$a^{o(1)}$ admissible points (Lemma~\ref{lem:fiber}). The number of triples $(u,v,D)$ is
$\ll (B/T)^2\cdot a/T^3$, where $B=a/A$, which gives~\eqref{eq:intro2}.

\subsection*{Formal verification}
Every step of the proof, including Theorem~\ref{thm:main} in the form stated above, has been
formalised in Lean~4~\cite{Lean} using Mathlib~\cite{Mathlib}. Section~\ref{sec:lean}
describes the formalisation. The constants in this paper are the constants in the Lean proof.

\subsection*{What is not proved}
Theorem~\ref{thm:main} is far from~\eqref{eq:conj}, and the method cannot reach it; see
Section~\ref{sec:limits}.

\section{Preliminaries}

\subsection{Notation}
Throughout, $a>b>0$ are integers, $b_1=b$, $b_{i+1}=a\bmod b_i$, and $P=P(a,b)$, so that
$a>b_1>b_2>\dots>b_P>0=b_{P+1}$. Write $q_i=\lfloor a/b_i\rfloor$, so
\[
a=q_ib_i+b_{i+1},\qquad 0\le b_{i+1}<b_i\qquad(1\le i\le P).
\]
Let $\tau(n)$ be the number of positive divisors of $n$, and let
\[
\Delta=\Delta(a)=\max_{1\le n\le 32a^7}\tau(n).
\]

\begin{lemma}[Quotients increase]\label{lem:quot}
If $b_{i+1}>0$ then $q_{i+1}\ge q_i+1$.
\end{lemma}

\begin{proof}
Since $b_{i+1}<b_i$ we have $(q_i+1)b_{i+1}\le q_ib_i+b_{i+1}=a$, so
$q_{i+1}=\lfloor a/b_{i+1}\rfloor\ge q_i+1$.
\end{proof}

\begin{lemma}[Divisor bound]\label{lem:div}
For every integer $k\ge1$ there is $C_k$ with $\tau(n)^k\le C_k\,n$ for all $n\ge1$.
\end{lemma}

\begin{proof}
This is standard: $\tau(n)^k/n=\prod_{p^\alpha\| n}(\alpha+1)^k/p^\alpha$. Each factor with
$p\ge 2^k$ is at most $(\alpha+1)^k/2^{k\alpha}\le1$. Each of the fewer than $2^k$ factors with
$p<2^k$ is at most $\sup_\alpha(\alpha+1)^k/2^\alpha<\infty$.
\end{proof}

\subsection{Window runs}

\begin{definition}
Let $A\ge1$. A \emph{window run} of length $T$ is a sequence $x_0,x_1,\dots,x_{T-1}$ of
consecutive terms of the sequence $(b_i)$, so $x_{j+1}=a\bmod x_j$, all lying in $[A,2A)$.
\end{definition}

In a window run, $x_0>x_1>\dots>x_{T-1}$ and $x_j\le a$.

\section{The curvature of three consecutive steps}

Fix four consecutive terms
\[
x=b_i,\quad y=b_{i+1}=a\bmod x,\quad z=b_{i+2}=a\bmod y,\quad w=b_{i+3}=a\bmod z,
\]
with $x,y,z>0$ (and $w\ge0$). Write
\[
q=\left\lfloor \frac ax\right\rfloor,\quad
u=\left\lfloor \frac ay\right\rfloor-q,\quad
v=\left\lfloor \frac az\right\rfloor-\left\lfloor \frac ay\right\rfloor,\qquad
d=x-y,\quad e=y-z,\quad f=z-w .
\]
By Lemma~\ref{lem:quot}, $u,v\ge1$, and clearly $d,e,f\ge1$. The division equations read
\begin{equation}\label{eq:div}
xq+y=a,\qquad y(q+u)+z=a,\qquad z(q+u+v)+w=a .
\end{equation}

\begin{definition}
The \emph{curvature} of the block is $D=dv-eu$.
\end{definition}

\begin{lemma}[Curvature identity]\label{lem:curv}
\[
D\,z=df+ude+(u-1)e^2 .
\]
In particular $D\ge1$.
\end{lemma}

\begin{proof}
Subtracting the first two equations of~\eqref{eq:div} gives $uy=qd+e$, hence
$uz=uy-ue=qd+(1-u)e$. Subtracting the last two gives $vz=(q+u)e+f$. Therefore
\[
(dv-eu)z=d\bigl((q+u)e+f\bigr)-e\bigl(qd+(1-u)e\bigr)=df+ude+(u-1)e^2 .
\]
The right-hand side is at least $df\ge1$, and $z>0$, so $D>0$.
\end{proof}

Equivalently, the identity is the linear combination
\[
Dz-\bigl(df+ude+(u-1)e^2\bigr)=d\,E_3+(a-qx-x-y+z)\,E_2+(-a+(q+u)y+y)\,E_1
\]
of the three division equations $E_1,E_2,E_3$ of~\eqref{eq:div} (each written as
$\text{left}-\text{right}=0$). This is how it is checked in Lean.

\begin{lemma}[The fiber equation]\label{lem:fiber-eq}
Put $M(q)=uq^2+(u^2+u+v)q+v$ and $N(q)=uq+u^2+v$. Then
\[
x\,M(q)=a\,N(q)+D .
\]
\end{lemma}

\begin{proof}
From~\eqref{eq:div}, $d=(q+1)x-a$ and $e=y-z=(q+u+1)y-a=a(q+u)-q(q+u+1)x$. Substituting
these into $D=dv-eu$ gives
$D=x\bigl(uq(q+u+1)+v(q+1)\bigr)-a\bigl(u(q+u)+v\bigr)$, which is the claim.
\end{proof}

\begin{lemma}[Size of the parameters]\label{lem:size}
If $x\le a$ then
\[
u\,yx\le a(d+e),\qquad v\,zy\le a(e+f),\qquad D\,z\le (u+1)(d+e+f)^2 .
\]
\end{lemma}

\begin{proof}
From $uy=qd+e$ we get $uyx=(qx)d+ex\le ad+ae$, since $qx\le a$ and $x\le a$. Similarly
$vz=(q+u)e+f$ with $(q+u)y\le a$ and $y\le a$ gives the second bound. For the third, by
Lemma~\ref{lem:curv},
$Dz=df+ude+(u-1)e^2\le df+u\,e(d+e)\le (1+u)(d+e+f)^2$.
\end{proof}

\section{Few blocks share a triple $(u,v,D)$}

\begin{lemma}[Fiber count]\label{lem:fiber}
Let $a\ge1$, $1\le u,v\le a$ and $1\le D\le 2a$. Let $M(q)=uq^2+pq+v$ with $p=u^2+u+v$, and
let $L(q)=cq+\ell$ with $c=au$ and $\ell=a(u^2+v)+D$. Then the set
\[
Q=\{q\ge 0:\ xM(q)=L(q)\text{ for some integer }x\ge1\}
\]
has at most $\Delta=\max_{n\le 32a^7}\tau(n)$ elements.
\end{lemma}

\begin{proof}
Put $U(q)=ucq+cp-u\ell$ and $R=u\ell^2-cp\ell+c^2v$. Expanding shows
\begin{equation}\label{eq:resultant}
c^2M(q)-U(q)L(q)=R
\end{equation}
identically in $q$. If $q\in Q$ with witness $x$, substituting $L(q)=xM(q)$
into~\eqref{eq:resultant} gives
\begin{equation}\label{eq:R}
R=M(q)\bigl(c^2-U(q)x\bigr).
\end{equation}
We bound the sizes first. From $u,v\le a$ and $D\le 2a$ we get $\ell\le 4a^3$, $p\le 3a^2$
and $c\le a^2$, so
\[
|R|\le u\ell^2+cp\ell+c^2v\le 16a^7+12a^7+a^5\le 32a^7,\qquad c^2\le a^4\le 32a^7 .
\]

\emph{Case $R\ne0$.} By~\eqref{eq:R}, $M(q)$ is a positive divisor of $|R|$. Since $u,p\ge1$,
the map $q\mapsto M(q)$ is strictly increasing on $q\ge0$, hence injective. Therefore
$|Q|\le\tau(|R|)\le\Delta$.

\emph{Case $R=0$.} Since $M(q)>0$, \eqref{eq:R} gives $U(q)x=c^2$. Here $c^2>0$ and $x\ge1$,
so $U(q)$ is a positive divisor of $c^2$. Since $U$ is linear with positive leading
coefficient $uc$, it is injective. Therefore $|Q|\le\tau(c^2)\le\Delta$.
\end{proof}

The case $R=0$ cannot be dropped: it is the case in which $M$ and $L$ might share a factor.

\section{Counting visits to a window}

In this section $x_0,\dots,x_{T-1}$ is a window run in $[A,2A)$, and
$\Delta=\max_{n\le 32a^7}\tau(n)$.

\begin{proposition}[Drop bound; after Erd\H{o}s--Shallit]\label{prop:drop}
If $T\ge2$ then $T(T-1)\le 4A\Delta$.
\end{proposition}

\begin{proof}
Let $h_j=x_j-x_{j+1}\ge1$ for $0\le j\le T-2$. By telescoping,
$\sum_j h_j=x_0-x_{T-1}\le (2A-1)-A=A-1$. Call $j$ \emph{good} if $h_j(T-1)\le 2A$ and
\emph{bad} otherwise. Each bad $j$ has $h_j(T-1)\ge 2A+1$, so
\[
|\Bad|\,(2A+1)\le (T-1)\textstyle\sum_j h_j\le (T-1)(A-1).
\]
If $2|\Bad|\ge T-1$ this would give $(T-1)(2A+1)\le 2(T-1)(A-1)$, which is false. Hence
$2|\Bad|<T-1$, and since $|\Good|+|\Bad|=T-1$ we get $2|\Good|\ge T$.

For good $j$ we have $1\le h_j\le H:=\lfloor 2A/(T-1)\rfloor$. Also
$a+h_j=q_jx_j+x_{j+1}+x_j-x_{j+1}=(q_j+1)x_j$, so $x_j\mid a+h_j$. Moreover
$a+h_j\le 2a\le 32a^7$. Since $j\mapsto x_j$ is injective, at most $\tau(a+h)\le\Delta$ good
indices have $h_j=h$. Hence $|\Good|\le H\Delta$, and
\[
T(T-1)\le 2|\Good|(T-1)\le 2\Delta\,H(T-1)\le 4A\Delta. \qedhere
\]
\end{proof}

\begin{proposition}[Curvature bound]\label{prop:curv}
If $T\ge12$ then $T^6A^2\le 995328\,a^3\Delta$.
\end{proposition}

\begin{proof}
For $0\le j\le T-4$ let $s_j=x_j-x_{j+3}$ be the total drop of the block
$x_j,x_{j+1},x_{j+2},x_{j+3}$. By telescoping,
\[
\sum_{j=0}^{T-4}s_j=(x_0+x_1+x_2)-(x_{T-3}+x_{T-2}+x_{T-1})\le 3(2A-1)-3A=3A-3 .
\]
Call $j$ \emph{good} if $s_jT\le 12A$ and \emph{bad} otherwise. Then
$|\Bad|(12A+1)\le T(3A-3)$. If $4|\Bad|\ge T$ this would give $T(12A+1)\le 4T(3A-3)$, which is
false. So $4|\Bad|\le T-1$. Since $|\Good|=T-3-|\Bad|$ and $T\ge12$,
\[
2|\Good|\ge 2T-6-\tfrac{T-1}2\ge T .
\]

Fix a good $j$. Apply Section~3 to $x=x_j$, so that $y,z,w=x_{j+1},x_{j+2},x_{j+3}$, all
in $[A,2A)$. Write $s=s_j=d+e+f$. Lemma~\ref{lem:size} gives $uA^2\le uyx\le a(d+e)\le as$.
Multiplying by $T$ and using $sT\le 12A$ gives $uA^2T\le 12aA$, that is,
\[
u\,AT\le 12a,\qquad\text{and likewise}\qquad v\,AT\le 12a .
\]
Also $DA\le Dz\le(u+1)s^2\le 2us^2$. Multiplying by $T^2$ gives
$DAT^2\le 2u(sT)^2\le 288uA^2$, so $DT^2\le 288uA$, and then
\[
D\,T^3\le 288\,u\,AT\le 3456\,a .
\]
Thus the triple $(u,v,D)$ of a good block lies in $[1,U]^2\times[1,D^*]$, where
\[
U=\lfloor 12a/(AT)\rfloor\le a,\qquad D^*=\lfloor 3456a/T^3\rfloor\le 2a .
\]
The inequalities $U\le a$ and $D^*\le 2a$ hold because $AT\ge12$ and $T^3\ge1728$.

Now fix such a triple and consider the good blocks that have it. By Lemma~\ref{lem:fiber-eq},
each one gives a point $q=q_j=\lfloor a/x_j\rfloor$ of the set $Q$ of Lemma~\ref{lem:fiber}.
The map $j\mapsto q_j$ is injective on these blocks: $q_j$ determines $x_j=L(q_j)/M(q_j)$, and
$x_j$ determines $j$. By Lemma~\ref{lem:fiber}, at most $\Delta$ good blocks share a triple.
So $|\Good|\le U^2D^*\Delta$, and
\[
T^6A^2=T\,(AT)^2\,T^3\le 2\Delta\,(U\,AT)^2\,(D^*T^3)\le 2\Delta(12a)^2(3456a)=995328\,a^3\Delta.
\qedhere
\]
\end{proof}

\begin{corollary}[Window bound]\label{cor:window}
If $T\ge12$ then $T^{10}\le 63700992\,a^3\Delta^3$.
\end{corollary}

\begin{proof}
Using $T\le 2(T-1)$ and the two propositions,
\[
T^{10}A^2\le 4\,(T^6A^2)\,\bigl(T(T-1)\bigr)^2\le 4\cdot 995328\,a^3\Delta\cdot(4A\Delta)^2
=63700992\,a^3\Delta^3A^2 .
\]
Divide by $A^2$.
\end{proof}

\begin{remark}
In terms of $B=a/A$, Propositions~\ref{prop:drop} and~\ref{prop:curv} say
$T\ll a^{1/2+o(1)}B^{-1/2}$ and $T\ll a^{1/6+o(1)}B^{1/3}$. The first is better when
$B>a^{2/5}$, and both equal $a^{3/10+o(1)}$ at $B=a^{2/5}$. Corollary~\ref{cor:window} takes
the product of the two bounds, so no case split is needed.
\end{remark}

\section{Proof of Theorem~\ref{thm:main}}

Let $0<b<a$. For $0\le k\le\log_2 a$ let $T_k$ be the number of indices $i\le P$ with
$2^k\le b_i<2^{k+1}$. Every $b_i$ with $i\le P$ lies in $[1,a)$, so
$P=\sum_{0\le k\le \log_2 a}T_k$. Since $(b_i)$ is decreasing, the indices counted by $T_k$
are consecutive. So when $T_k>0$ they form a window run in $[A,2A)$ with $A=2^k$.

Let $\varepsilon>0$. Choose an integer $m>7/(5\varepsilon)$, and let $C_{3m}$ be the constant of
Lemma~\ref{lem:div} with $k=3m$. The maximum $\Delta$ is attained at some $n_0\le 32a^7$, so
$\Delta^{3m}\le C_{3m}\cdot 32a^7$. If $T_k\ge12$, Corollary~\ref{cor:window} gives
\[
T_k^{10m}\le 63700992^m\,a^{3m}\,\Delta^{3m}\le K\,a^{3m+7},\qquad K=63700992^m\cdot C_{3m}\cdot 32,
\]
and hence
\[
T_k\le K^{1/(10m)}\,a^{3/10+7/(10m)}\le K^{1/(10m)}\,a^{3/10+\varepsilon/2}.
\]
If $T_k\le11$ this bound holds with $11$ in place of $K^{1/(10m)}$. Finally
$\lfloor\log_2a\rfloor+1\le C'_\varepsilon a^{\varepsilon/2}$. Summing over $k$,
\[
P(a,b)\le (\lfloor\log_2 a\rfloor+1)\bigl(11+K^{1/(10m)}\bigr)a^{3/10+\varepsilon/2}
\le C_\varepsilon\, a^{3/10+\varepsilon}. \qquad\qed
\]

\section{The Lean formalisation}\label{sec:lean}

The repository accompanying this paper contains a Lean~4 project (toolchain
\lean{v4.35.0}, Mathlib \lean{v4.35.0}) that proves Theorem~\ref{thm:main}. The final
statement is
\begin{quote}\ttfamily\small
theorem P\_le : $\forall$ $\varepsilon$ > (0 : $\mathbb R$), $\exists$ C > 0, $\forall$ a b : $\mathbb N$, 0 < b $\to$ b < a $\to$\\
\hspace*{2em}(P a b : $\mathbb R$) $\le$ C * (a : $\mathbb R$) \^{} ((3 : $\mathbb R$) / 10 + $\varepsilon$)
\end{quote}
where \lean{P a b} is defined by well-founded recursion as the number of positive terms
($P(a,0)=0$, $P(a,b)=P(a,a\bmod b)+1$ for $b>0$), and the file checks $P(35,22)=7$.
The command \lean{\#print axioms P\_le} reports only \lean{propext}, \lean{Classical.choice}
and \lean{Quot.sound}. There is no \lean{sorry} and no \lean{native\_decide}.

\begin{center}\small
\begin{tabular}{@{}lll@{}}
\toprule
Paper & Lean & File \\
\midrule
Definition of $P(a,b)$ & \lean{P}, \lean{orb} & \lean{Defs.lean}\\
Lemma~\ref{lem:quot} & \lean{div\_lt\_div\_mod} & \lean{Block.lean}\\
Lemma~\ref{lem:curv} & \lean{curvature\_identity}, \lean{one\_le\_bD} & \lean{Block.lean}\\
Lemma~\ref{lem:fiber-eq} & \lean{fiber\_identity}, \lean{block\_fiber\_eq} & \lean{Block.lean}\\
Lemma~\ref{lem:size} & \lean{bu\_mul\_le}, \lean{bv\_mul\_le}, \lean{bD\_mul\_le} & \lean{Block.lean}\\
Lemma~\ref{lem:fiber} & \lean{fiber\_poly\_identity}, \lean{fiber\_card\_le} & \lean{Fiber.lean}\\
Window runs & \lean{Window}, \lean{DivBound} & \lean{Window.lean}\\
Proposition~\ref{prop:drop} & \lean{drop\_bound} & \lean{Window.lean}\\
Proposition~\ref{prop:curv} & \lean{curvature\_bound} & \lean{Window.lean}\\
Corollary~\ref{cor:window} & \lean{window\_bound} & \lean{Window.lean}\\
Dyadic decomposition & \lean{P\_eq\_sum\_visits}, \lean{visits\_window} & \lean{Main.lean}\\
Theorem~\ref{thm:main} & \lean{P\_le} & \lean{Main.lean}\\
Lemma~\ref{lem:div} & \lean{exists\_divisors\_power\_bound} & \lean{DivisorPowerBound.lean}\\
\bottomrule
\end{tabular}
\end{center}

Some details of the formalisation:
\begin{itemize}
\item The algebraic identities (Lemmas~\ref{lem:curv}, \ref{lem:fiber-eq}
and~\eqref{eq:resultant}) are proved over $\mathbb Z$ by the tactic
\lean{linear\_combination}, using the explicit coefficients above. They are then transferred to $\mathbb N$ by writing
$z=w+f$, $y=z+e$, $x=y+d$, which avoids truncated subtraction.
\item The curvature is the natural number $\lean{bD}\,a\,x=(x-y)v-(y-z)u$. Lean first proves
$eu<dv$, so this truncated subtraction is the true value.
\item The counting arguments use \lean{Finset.card\_le\_mul\_card\_image\_of\_maps\_to}: the
number of good indices is at most the size of the parameter box times the largest fiber.
\item The divisor bound $\tau(n)^k\le C_kn$ (Lemma~\ref{lem:div}) is the file
\lean{DivisorPowerBound.lean} from the \lean{openai/math} repository~\cite{OAIMath}
(Apache-2.0), vendored unchanged except for its import lines.
\item The project has about 1000 lines of Lean. \lean{lake build} rebuilds it in under a
minute once Mathlib is built.
\end{itemize}

\section{Limits of the method}\label{sec:limits}

The argument counts the triples $(u,v,D)$ independently over their ranges. Even if the divisor
loss in Lemma~\ref{lem:fiber} were replaced by an absolute constant, the bound would still be
$T\ll(B/T)^2\,a/T^3$. Combined with $T\ll A^{1/2}$, this still gives the exponent $3/10$. So
removing the $\varepsilon$, sharpening constants or improving the divisor count cannot give
a logarithmic bound. Further progress needs to restrict \emph{which} triples can occur together
along a single orbit.

There is no obvious monotonicity to exploit. The curvatures $D_i$ of successive blocks need
not decrease. For $a=13242$ the orbit of $255$ begins
\[
255,\ 237,\ 207,\ 201,\ 177,\ 144,\ 138,\ 132,\ \dots
\]
with quotients $51,55,63,65,74,91,95,100,\dots$. The curvatures of the first five blocks are
$24,12,6,111,30$, all for blocks that lie within a factor of two.

\subsection*{Acknowledgements}
The argument was developed with the help of AI language models. All proofs were then verified
in Lean~4 against Mathlib. We thank the Mathlib community, and the authors of \lean{openai/math}
for the formalised divisor bound.

\begin{thebibliography}{99}

\bibitem{CP}
Z.~Chase and M.~Pandey,
\emph{On the length of Pierce expansions},
arXiv:2211.08374 (2022).

\bibitem{ES}
P.~Erd\H{o}s and J.~O.~Shallit,
\emph{New bounds on the length of finite Pierce and Engel series},
S\'em.\ Th\'eor.\ Nombres Bordeaux (2) \textbf{3} (1991), 43--53.

\bibitem{Lean}
L.~de~Moura and S.~Ullrich,
\emph{The Lean~4 theorem prover and programming language},
in: Automated Deduction -- CADE 28, Lecture Notes in Comput.\ Sci.\ \textbf{12699},
Springer, 2021, 625--635.

\bibitem{Mathlib}
The mathlib Community,
\emph{The Lean mathematical library},
in: Proceedings of the 9th ACM SIGPLAN International Conference on Certified Programs and
Proofs (CPP 2020), 367--381.

\bibitem{OAIMath}
OpenAI,
\texttt{openai/math}, \url{https://github.com/openai/math}.

\bibitem{OPG}
J.~O.~Shallit,
\emph{A discrete iteration related to Pierce expansions},
Open Problem Garden,
\url{https://www.openproblemgarden.org/op/a_discrete_iteration_related_to_pierce_expansions}.

\bibitem{Pierce}
T.~A.~Pierce,
\emph{On an algorithm and its use in approximating roots of algebraic equations},
Amer.\ Math.\ Monthly \textbf{36} (1929), 523--525.

\bibitem{Shallit86}
J.~O.~Shallit,
\emph{Metric theory of Pierce expansions},
Fibonacci Quart.\ \textbf{24} (1986), 22--40.

\bibitem{ShallitTalks}
J.~O.~Shallit,
\emph{Talks}, \url{https://cs.uwaterloo.ca/~shallit/talks.html}
(accessed October 2026).

\end{thebibliography}

\end{document}
